\chapter{Current Validation and Observations}
\label{chap:validation}

This chapter reports two different kinds of evidence, and keeps them
clearly apart. The first is an automated test suite that checks each
component's logic in isolation, against mocked external systems. The
second is a small, real-system pilot that checks whether the same
components behave correctly against the real PyPI registry, a real Ollama
model, a real Docker daemon, and real GitHub repositories. Neither of these
is the thesis's final evaluation, which requires running the full 187-row
evaluation split described in \Cref{chap:problem-dataset} and has not yet
been performed (\Cref{chap:future-work}).

\section{Automated Testing}
\label{sec:v-testing}

At the implementation state this thesis reports, the repository's test
suite contains \textbf{310 passing tests}, organised by the component they
exercise.

\begin{table}[h]
  \centering
  \begin{tabular}{|p{5.2cm}|r|p{5.8cm}|}
    \hline
    \textbf{Area} & \textbf{Tests} & \textbf{What is checked} \\
    \hline
    Context extraction and classification (I2) & 9 &
      Root-cause hints, confidence levels, and the legacy-hint
      acceptance/rejection logic from \Cref{sec:m-classification}. \\
    \hline
    Explanation schema, prompt rendering, and runner (I3) & 29 &
      Prompt template rendering, JSON schema validation, and the
      retry-on-failure logic from \Cref{sec:m-explanation}. \\
    \hline
    PyPI retrieval and compatibility evidence (I4) & 120 &
      Import-name mapping, candidate-version filtering, Python-compatibility
      checks, and the compatibility-evidence registry lookup from
      \Cref{sec:m-repair-retrieval,sec:m-repair-compat}. \\
    \hline
    Repair proposal validation and agent orchestration (I4) & 78 &
      Grounding validation, argument-list construction, and rejection of
      an invented name or an out-of-range version, from \Cref{sec:m-repair}. \\
    \hline
    Fix application, Docker orchestration, and outcome classification (I5) & 74 &
      Command re-validation, the dataset re-join, outcome classification,
      and cleanup behaviour, from \Cref{sec:m-fixapply}. \\
    \hline
    \textbf{Total} & \textbf{310} & \\
    \hline
  \end{tabular}
  \caption{Automated tests by component area, at the I5 implementation state.}
  \label{tab:v-tests}
\end{table}

Dataset preparation (I1) has no dedicated automated test file of its own;
instead, its output is checked by independently recomputing the counts in
\Cref{chap:problem-dataset} directly from its output files and confirming
they agree exactly, which they do.

Across these areas, the tests specifically exercise: rejection of malformed
or schema-invalid model output; rejection of a repair proposal that invents
a package name or a version outside the retrieved evidence; the shell-safety
checks described in \Cref{sec:m-fixapply}, including a static check that
neither \texttt{shell=True}, \texttt{os.system}, \texttt{eval}, nor
\texttt{exec} appears anywhere in the components that build or run a
command; the three-outcome classification logic, including the conservative
same-error check; and cleanup behaviour under both success and failure. All
of this runs against mocked subprocesses, a mocked Ollama endpoint, and
mocked HTTP responses. This is deliberate and also a limitation: a mocked
test can prove a component's logic is internally consistent, but, as
\Cref{sec:m-fixapply-bugs} already showed for three real bugs it could not
have caught, it cannot prove the component behaves correctly against a real
operating system, a real Docker daemon, or a real network. That is exactly
what the pilots below were run to check.

\section{I4 Implementation Validation}
\label{sec:v-i4}

\subsection{Retrieval Against Real PyPI}

The five import names used in the original L5 design check
(\texttt{sklearn}, \texttt{umap}, \texttt{pkg\_resources}, \texttt{scipy},
\texttt{dms\_variants}) were re-checked against the real, production
\texttt{retrieve()} entry point and the live PyPI registry. All four mapped
names resolved correctly to their real PyPI distributions with plausible,
current candidate versions; \texttt{dms\_variants}, which has no verified
mapping, correctly returned \texttt{mapping\_unknown} with no PyPI request
made at all, confirming under real conditions -- not only in a mocked test
-- that an unmapped import never reaches the network.

A later, separate check confirmed a fix to the mapping table itself.
\texttt{cv2} and \texttt{skimage} are both common, usable failing modules in
the dataset (\Cref{tab:pd-top-modules}) that had no entry in
\texttt{config/package\_mapping.yaml}, so the retriever returned
\texttt{mapping\_unknown} for both instead of resolving them. Adding
\texttt{cv2: opencv-python} and \texttt{skimage: scikit-image} was checked
against the real PyPI registry, alongside a regression check confirming
that the existing \texttt{sklearn} mapping still resolved correctly and was
not disturbed by the addition.

\subsection{Repair Proposals Against a Real Model}

A small pilot ran \texttt{rag\_repair\_agent.py}'s real entry point, using a
real local Ollama model, over four records:

\begin{itemize}
  \item An excluded, \texttt{system\_library} record correctly abstained
        with zero PyPI or Ollama calls, confirming the eligibility gate
        works under real conditions.
  \item A \texttt{missing\_package} record (\texttt{sklearn}) timed out
        twice under the originally configured 120-second ceiling, then
        succeeded on the first attempt once the ceiling was raised to 300
        seconds -- the direct evidence behind the timeout change reported
        in \Cref{sec:m-repair}.
  \item The same record, rerun at 300 seconds, produced a schema-valid,
        grounded proposal: install \texttt{scikit-learn}.
  \item A \texttt{wrong\_version} record (the \texttt{scipy.integrate.cumtrapz}
        case from \Cref{tab:m-compat}) produced a grounded
        \texttt{pin\_version} proposal for \texttt{scipy==1.13.1}, using the
        real compatibility-evidence intersection on the first attempt.
\end{itemize}

\subsection{What This Pilot Does and Does Not Show}

It shows that real PyPI retrieval, real import-name mapping, and real model
proposal generation work end to end, and that timeout and retry behaviour
degrades safely rather than hanging or fabricating a result. It does not
show an installed package, a modified environment, or a re-executed
notebook -- none of that happened at this step, by design; validating a
proposal is FixApplicator's job, covered next.

\section{I5 Implementation Validation}
\label{sec:v-i5}

A second pilot ran \texttt{apply\_and\_validate()} against a real Docker
daemon and real GitHub repositories, reusing already-validated I4 proposals.
The upstream sibling database was found locked by another process during
this pilot, which meant commit metadata for one lookup had to fall back to
its documented "unavailable" behaviour -- itself a real, unplanned exercise
of that fallback path, not a simulated one.

\begin{table}[h]
  \centering
  \small
  \begin{tabular}{|p{1.7cm}|p{3.2cm}|p{2.6cm}|p{6.0cm}|}
    \hline
    \textbf{Case} & \textbf{Original error} & \textbf{Fix applied} & \textbf{Result} \\
    \hline
    1 -- \texttt{sklearn} & \texttt{ModuleNot\-FoundError}: sklearn & install \texttt{scikit-learn} &
      Original error gone; \texttt{still\_failing} on a later, unrelated \texttt{No module named 'pandas'}. \\
    \hline
    2 -- \texttt{scipy} & \texttt{ImportError}: cumtrapz & pin \texttt{scipy==1.13.1} &
      Original error gone; \texttt{still\_failing} on a later, unrelated \texttt{No module named 'tf\_keras'}. \\
    \hline
    3 -- controlled failure & synthetic: non-existent repository & n/a &
      \texttt{apply\_error}, failure stage \texttt{clone}; no Docker artefacts created. \\
    \hline
    4 -- \texttt{cv2} & \texttt{ModuleNot\-FoundError}: cv2 & install \texttt{opencv-python} &
      Original error gone; \texttt{still\_failing} on a missing \emph{system} library (\texttt{libxcb.so.1}), a known characteristic of \texttt{opencv-python} on a minimal base image. \\
    \hline
    5 -- \texttt{sklearn} (2nd repo) & \texttt{ModuleNot\-FoundError}: sklearn & install \texttt{scikit-learn} &
      Original error gone; \texttt{still\_failing} on a later, unrelated \texttt{No module named 'tensorflow'}. \\
    \hline
  \end{tabular}
  \caption{I5 real-Docker pilot cases. Case 3 is a deliberately synthetic
  infrastructure-failure test, not a dataset row.}
  \label{tab:v-i5-cases}
\end{table}

Cases 4 and 5 were chosen deliberately, after cases 1--3, specifically to
give the pilot the best possible chance of reaching a completely clean
\texttt{fixed} result: both repositories have no declared requirements
file, no \texttt{setup.py}, and exactly one notebook, minimising the chance
of a second, unrelated dependency gap.

\subsection{Key Observation: Fixing One Dependency Can Reveal Another}
\label{sec:v-key-observation}

Every one of the five real cases above -- spanning four different
repositories, two different fix actions, and packages with no shared code
lineage -- resolved the specific error it targeted and then failed again on
a different, previously invisible problem. This is not a coincidence and it
is not a defect in FixApplicator's own classification logic, which is
separately unit-tested and proven capable of reporting a clean
\texttt{fixed} result the moment one occurs. It is a property of the
dataset: the Docker pipeline's error classification records only the
\emph{first} import failure a notebook hits. A real research notebook can
easily depend on several packages that are each individually missing; once
the first one is repaired, execution proceeds further and can reach the
next one, which was simply invisible until that point. The conservative
\texttt{same\_as\_original\_error} check (\Cref{chap:methodology})
correctly distinguished "a different, new problem" from "the fix did not
work" in every one of these cases.

No completely clean \texttt{fixed} outcome has yet been produced, despite a
pilot deliberately designed to find one. This is stated plainly rather than
minimised: a clean outcome remains achievable in principle on a
single-dependency notebook, or with the future two-round repair mode
feeding a \texttt{still\_failing} result's new error back into a second
repair attempt (\Cref{chap:future-work}), but neither has been demonstrated
yet.

\section{Summary: Implementation Validation, Not Final Evaluation}
\label{sec:v-summary}

The results in this chapter show that every implemented component -- error
classification, explanation, PyPI-grounded repair proposal generation, and
fix application with re-execution -- works correctly end to end against
real external systems, not only against mocked fixtures. They also
surfaced concrete, fixable problems (a timeout set too low, two missing
package mappings, three environment-specific bugs) that only a real pilot
could have found.

What these results do \emph{not} show is an overall repair success rate.
Five real cases, or four PyPI names, or four repair proposals is not a
sample large enough to support a general claim about how often this system
repairs a dependency error, and this thesis makes no such claim. That
question requires running the full 187-row evaluation split
(\Cref{chap:problem-dataset}) under a controlled, orchestrated run, which is
described as remaining work in \Cref{chap:future-work}.
